\chapter{A wick Tutorial}

This chapter builds a small but complete game, from an empty folder to
something you can play. We introduce language features as we need them and
explain them only as far as the game requires; the precise rules wait for
Chapter~3. By the end you will have written variables, functions, an optional,
a list, a loop, and the two functions the engine calls every frame --- and you
will have used the two tools that make wick pleasant to work in, hot reload
and the on-screen error screen.

\section{Installing and running}

wick ships inside the lantern engine. There is nothing separate to install:
if you can build the engine, you can run wick.

\begin{lstlisting}[language=bash,basicstyle=\ttfamily\small,keywordstyle={}]
git clone https://github.com/alikatgh/lantern && cd lantern
brew install sdl2 lua cmake pkg-config   # macOS; Linux: same packages
cmake -B build && cmake --build build -j8

./build/lantern games/wicklab            # the first wick program
./build/lantern games/showcase_wick      # Kora Night, in wick
\end{lstlisting}

\noindent
The first command builds the engine. The last two run games that already
exist in the repository: `wicklab`, a small tour of the language, and
`showcase_wick`, the wick build of Kora Night that recurs throughout this
book.

\section{A game is a folder}

A wick game is a directory that contains a file named `main.wick`. When you
point the engine at a folder, it runs that file if present, falling back to
`main.lua` otherwise. To start a new game:

\begin{lstlisting}[language=bash,basicstyle=\ttfamily\small,keywordstyle={}]
mkdir games/lamp
$EDITOR games/lamp/main.wick
./build/lantern games/lamp
\end{lstlisting}

\noindent
The screen is 400 by 240 pixels. The engine runs at 60 frames per second and
looks up two functions by name each frame:

\begin{lstlisting}[language=wick]
fn update(dt: num) {   // dt = seconds since last frame, capped at 0.1
}

fn draw() {            // 3D first, then 2D composites on top
  lt.clear(0.1, 0.1, 0.2)
  lt.print("HELLO 400X240", 4, 4, 1, 1, 1, 1)
}
\end{lstlisting}

\noindent
Both functions are optional; a `main.wick` with only top-level statements is
a legal program. Top-level statements --- everything not inside a `fn` --- run
once, at load, and that is where you create resources and set up state. Save
the file with those two functions in it, run the folder, and you have a
window that clears to a dark blue and prints a line of text in the built-in
8-by-8 font. That is the smallest thing worth calling a game.

\section{Variables and the first motion}

Let us make something move. Every variable is introduced with `let`. Types
are inferred from the initialiser, so you rarely write them:

\begin{lstlisting}[language=wick]
let x = 200.0      // num, inferred
let y = 120.0
\end{lstlisting}

\noindent
Those are top-level `let`s, so `x` and `y` are program globals: they persist
across frames, which is exactly what we want for the player's position. Now
read the keyboard in `update` and draw a square at \texttt{(x, y)} in `draw`:

\begin{lstlisting}[language=wick]
let x = 200.0
let y = 120.0

fn update(dt: num) {
  let speed = 120.0
  if lt.key("left")  { x = x - speed * dt }
  if lt.key("right") { x = x + speed * dt }
  if lt.key("up")    { y = y - speed * dt }
  if lt.key("down")  { y = y + speed * dt }
  x = max(8, min(392, x))
  y = max(8, min(232, y))
}

fn draw() {
  lt.clear(0.09, 0.07, 0.16)
  lt.rect(x - 4, y - 4, 8, 8, 1, 0.85, 0.3, 1)
}
\end{lstlisting}

\noindent
A few things are already worth noticing. `speed` is declared with `let`
inside `update`, so it is a local: it exists only for that call. Multiplying
by `dt`, the seconds elapsed since the last frame, makes the motion
frame-rate independent --- the square moves 120 pixels per second whether the
machine runs at 60 or 144 frames. `max` and `min` are built-ins, always
available with no namespace. And `lt.key` returns a `bool`, which is the only
thing an `if` will accept: there is no truthiness in wick, so you cannot
write `if x` and hope it means ``if x is nonzero.'' It simply will not
compile.

\section{Functions}

A function is declared with `fn`, always at the top level --- functions do not
nest in v0.1. Parameter types are required; the return type follows the
parameter list after a colon, and you omit it for a function that returns
nothing. Here is a tiny helper that makes the square bob up and down:

\begin{lstlisting}[language=wick]
fn wobble(v: num): num {
  return sin(v * 2) * 6
}
\end{lstlisting}

\noindent
`sin` is a built-in and takes radians. To use `wobble` we need a clock, so
add a top-level `let t = 0.0` and advance it in `update` with `t = t + dt`.
Then the draw becomes:

\begin{lstlisting}[language=wick]
lt.rect(x - 4, y - 4 + wobble(t), 8, 8, 1, 0.85, 0.3, 1)
\end{lstlisting}

\noindent
One rule to remember: functions must be declared before they are called
(the compiler is single-pass; Chapter~7 explains why). `update` and `draw`
are called by the host, not by your code, so their position relative to each
other does not matter --- but `wobble` must appear above the line that calls
it.

\section{An optional, and a saved score}

Now the part that makes wick wick. We want to keep a running count and save
it to disk. Saving is easy:

\begin{lstlisting}[language=wick]
lt.save("lamp_best", str(best))
\end{lstlisting}

\noindent
`str` converts a `num` to a `str`; `lt.save` writes bytes to a named save
slot. Loading is where the type system earns its keep. `lt.load_save` cannot
promise a value --- the save may not exist, or may be unreadable --- so its
return type is not `str` but `str?`, an \emph{optional} string. You cannot
use a `str?` where a `str` is wanted; the compiler stops you:

\begin{lstlisting}[language=wick]
let s = lt.load_save("lamp_best")   // s: str?
lt.print(s, 4, 4, 1, 1, 1, 1)       // COMPILE ERROR: expected str, got str?
\end{lstlisting}

\noindent
The error is `expected str, got str?`, and it is the good kind of error: it
happens the moment you write the line, not on some player's machine a month
after release. To turn a `str?` into a `str` you must say what happens when
it is `nil`. The quickest way is `??`, which supplies a default:

\begin{lstlisting}[language=wick]
let best = num(lt.load_save("lamp_best") ?? "") ?? 0
\end{lstlisting}

\noindent
Read it inside out. `lt.load_save("lamp_best")` is a `str?`. The first `??`
turns it into a `str`, using `""` when the save is absent. `num` parses a
`str` into a number, but parsing can fail --- `num("banana")` and even
`num("")` cannot produce a number --- so `num` returns `num?`. The second `??`
turns that into a plain `num`, using `0` when the parse failed. Both ways the
value can be missing are handled, in one line, because the compiler refused
to let either slip through. This is the zero-byte-save bug from the
introduction, made unwritable. Chapter~4 tells the whole story.

Wire it up: declare `best` at the top level, and bump it when the player
presses \texttt{Z}:

\begin{lstlisting}[language=wick]
let best = num(lt.load_save("lamp_best") ?? "") ?? 0

fn update(dt: num) {
  // ... movement from before ...
  if lt.pressed("z") {
    best = best + 1
    lt.save("lamp_best", str(best))
  }
}
\end{lstlisting}

\noindent
`lt.pressed` is true only on the frame a key goes down, where `lt.key` is
true for as long as it is held. Run the game, press \texttt{Z} a few times,
quit, and run it again: the count survives, because it was saved and safely
reloaded.

\section{A list: the lamps}

The game is called \emph{lamp}, so let us add lamps. We will keep four of
them, each with a position and a lit/unlit flag. wick has no structs in v0.1,
so we use the idiom that stands in for them: \emph{parallel lists}, where
index `i` names one lamp across several lists.

\begin{lstlisting}[language=wick]
let lamp_x: list<num> = []
let lamp_z: list<num> = []
let lamp_lit: list<bool> = []
\end{lstlisting}

\noindent
An empty list literal `[]` has no elements to infer a type from, so it
\emph{requires} an annotation --- that is what `list<num>` on the left is for.
Fill the lists with four lamps:

\begin{lstlisting}[language=wick]
push(lamp_x, 3.4)  push(lamp_z, 3.4)  push(lamp_lit, true)
push(lamp_x, -3.4) push(lamp_z, 3.4)  push(lamp_lit, true)
push(lamp_x, -3.4) push(lamp_z, -3.4) push(lamp_lit, true)
push(lamp_x, 3.4)  push(lamp_z, -3.4) push(lamp_lit, true)
\end{lstlisting}

\noindent
`push` appends, and it is type-checked: `push(lamp_lit, 3)` would not compile,
because `lamp_lit` holds `bool`. Lists are 0-based and you iterate them with a
numeric range. Here is a function that relights every lamp, using a `for` loop
over the half-open range \texttt{[0, 4)}:

\begin{lstlisting}[language=wick]
fn reset_lamps() {
  for i in 0..4 { lamp_lit[i] = true }
}
\end{lstlisting}

\noindent
`for i in a..b` runs `i` from `a` up to but not including `b` --- end-exclusive,
to match 0-based indexing, so `0..4` and `0..len(lamp_lit)` mean the same
thing here. Reading or writing an index that is out of range is a runtime
error with a line number, not a silent `nil`; we will meet that error screen
in a moment.

\section{Putting a round together}

We now have enough to make a game with a goal: lamps go dark over time, and
the player relights the nearest one by standing next to it and pressing
\texttt{Z}. The full logic is in `games/showcase_wick/main.wick`; here is the
shape of the update, condensed to the ideas we have covered.

\begin{lstlisting}[language=wick]
let score = 0

fn nearest_unlit(): num {     // lamp index, or -1 if all lit
  let bi = 0 - 1
  let bd = 0.0
  for i in 0..4 {
    if not lamp_lit[i] {
      let dx = lamp_x[i] - x
      let dz = lamp_z[i] - y
      let d = dx * dx + dz * dz
      if bi < 0 or d < bd { bi = i  bd = d }
    }
  }
  return bi
}

fn update(dt: num) {
  // ... movement, save on Z ...
  let i = nearest_unlit()
  if i >= 0 {
    if lt.pressed("z") { lamp_lit[i] = true  score = score + 1 }
  }
}
\end{lstlisting}

\noindent
Two things here look forward to later chapters. `nearest_unlit` returns a
plain `num` --- the index, or `-1` as a sentinel --- because wick v0.1 has no
multiple return values, so returning ``the index or nothing'' is done with an
out-of-band value rather than an optional. And the nested `if` on the last
two lines, rather than one `if i >= 0 and lt.pressed("z")`, keeps the
condition simple; Chapter~3 covers when the compiler asks you to parenthesise
combined conditions.

\section{The development loop}

Two features turn the write-run-fix cycle from minutes into seconds.

\paragraph{Hot reload.} Save `main.wick` while the engine is running and it
recompiles and reloads live. All previously created resources are freed
first, so reloading never leaks --- you can iterate on gameplay for an hour
without restarting. Change a constant, save, and watch the running game
change.

\paragraph{The error screen.} When your file has a mistake, the engine does
not crash to the terminal. It renders the error in-engine, as
\texttt{file:line: message}, and keeps waiting. Compile errors --- a `str?`
used as a `str`, a `+` between a string and a number, a call to a function
defined below its use --- and runtime errors alike --- an out-of-range index,
a failed `check()` --- all show the same way. Fix the file, save, and the game
hot-reloads on top of the error. You rarely leave the running game.

For example, deleting the `?? 0` from our `best` line produces, on screen:

\begin{lstlisting}[basicstyle=\ttfamily\small,keywordstyle={}]
main.wick:1: expected num, got num?
\end{lstlisting}

\noindent
and the game waits, showing you the line, until you put it back.

\paragraph{Deterministic screenshots.} For automated testing, point
`LANTERN_SHOT` at a path and the engine captures frame~60 as a BMP and then
exits. Add `LANTERN_FIXED_DT=1` and that capture is byte-identical on every machine,
because the frame clock is fixed-step and wick's `rand()` is deterministic.
Whole gameplay sessions replay exactly. Chapter~6 builds a testing discipline
on this; for now it is enough to know that ``it looks right on my machine''
can be turned into a check that means the same thing on every machine.

That is a complete game and the whole workflow around it. The rest of the
book makes precise what this chapter showed by example.
